\documentclass[11pt,a4paper]{article}

\usepackage[utf8]{inputenc}
\usepackage[T1]{fontenc}
\usepackage[spanish]{babel}
\usepackage{geometry}
\geometry{margin=2.5cm}

\usepackage{listings}
\usepackage{xcolor}
\usepackage{hyperref}
\usepackage{booktabs}
\usepackage{enumitem}
\usepackage{graphicx}
\usepackage{fancyhdr}
\usepackage{titlesec}
\usepackage{mdframed}
\usepackage{amsmath}

% Colores
\definecolor{bg}{RGB}{30,30,30}
\definecolor{fg}{RGB}{220,220,220}
\definecolor{kw}{RGB}{86,156,214}
\definecolor{cm}{RGB}{106,153,85}
\definecolor{st}{RGB}{206,145,120}
\definecolor{nb}{RGB}{78,201,176}
\definecolor{warnbg}{RGB}{255,248,220}
\definecolor{warnborder}{RGB}{200,150,0}
\definecolor{tipbg}{RGB}{220,240,255}
\definecolor{tipborder}{RGB}{0,100,200}

% Estilo listings — consola oscura
\lstdefinestyle{console}{
  backgroundcolor=\color{bg},
  basicstyle=\ttfamily\small\color{fg},
  keywordstyle=\color{kw}\bfseries,
  commentstyle=\color{cm}\itshape,
  stringstyle=\color{st},
  numberstyle=\tiny\color{gray},
  breaklines=true,
  breakatwhitespace=false,
  frame=single,
  framerule=0.5pt,
  rulecolor=\color{gray!50},
  xleftmargin=1em,
  xrightmargin=0.5em,
  aboveskip=0.8em,
  belowskip=0.8em,
  showstringspaces=false,
}

% Estilo listings — código macro Geant4
\lstdefinestyle{macro}{
  backgroundcolor=\color{bg},
  basicstyle=\ttfamily\small\color{fg},
  commentstyle=\color{cm}\itshape,
  morecomment=[l]{\#},
  keywordstyle=\color{kw},
  morekeywords={/detector,/ambe,/run,/vis,/event,/tracking,/control},
  stringstyle=\color{st},
  frame=single,
  framerule=0.5pt,
  rulecolor=\color{gray!50},
  xleftmargin=1em,
  xrightmargin=0.5em,
  aboveskip=0.8em,
  belowskip=0.8em,
  showstringspaces=false,
  breaklines=true,
}

% Cajas de advertencia y tip
\newmdenv[
  backgroundcolor=warnbg,
  linecolor=warnborder,
  linewidth=2pt,
  topline=false, bottomline=false,
  rightline=false,
  innerleftmargin=10pt,
  innerrightmargin=10pt,
  innertopmargin=6pt,
  innerbottommargin=6pt,
]{advertencia}

\newmdenv[
  backgroundcolor=tipbg,
  linecolor=tipborder,
  linewidth=2pt,
  topline=false, bottomline=false,
  rightline=false,
  innerleftmargin=10pt,
  innerrightmargin=10pt,
  innertopmargin=6pt,
  innerbottommargin=6pt,
]{tip}

% Encabezado / pie
\pagestyle{fancy}
\fancyhf{}
\rhead{\textcolor{gray}{\small Neutron Thermalization — AmBe}}
\lhead{\textcolor{gray}{\small Guía de visualización}}
\cfoot{\thepage}
\renewcommand{\headrulewidth}{0.4pt}

% Títulos con número en color
\titleformat{\section}{\large\bfseries\color{black}}{%
  \colorbox{black}{\textcolor{white}{\thesection}}}{0.6em}{}
\titleformat{\subsection}{\normalsize\bfseries}{%
  \thesubsection}{0.5em}{}

\hypersetup{
  colorlinks=true,
  linkcolor=blue!70!black,
  urlcolor=blue!70!black,
}

% -------------------------------------------------------------------
\begin{document}

\begin{titlepage}
  \centering
  \vspace*{3cm}
  {\Huge\bfseries Guía: de macro a video\par}
  \vspace{0.5cm}
  {\large Simulación de termalización de neutrones con fuente AmBe\par}
  \vspace{0.3cm}
  {\large Geant4 + Python + ffmpeg\par}
  \vspace{2cm}
  \rule{\textwidth}{1pt}
  \vspace{0.4cm}
  {\normalsize
    \begin{tabular}{ll}
      Proyecto:   & Neutron\_Thermalization \\
      Branch:     & \texttt{ambe\_source} \\
      Simulador:  & Geant4 11.04-beta-01 (OGLSX + gl2ps 1.4.2) \\
      Entorno Python: & \texttt{mi\_entorno} (virtualenv) \\
      Fecha:      & Mayo 2026 \\
    \end{tabular}
  }
  \vspace{0.5cm}
  \rule{\textwidth}{1pt}
  \vspace{2cm}

  \begin{tip}
    \textbf{Resumen del flujo:}
    \begin{enumerate}[noitemsep]
      \item Editar \texttt{vis\_video.mac} (geometría, fuente, colores).
      \item Correr Geant4 $\to$ genera 120 archivos EPS por evento.
      \item Convertir EPS $\to$ PNG con Ghostscript.
      \item Acumular frames con Python (\texttt{accumulate\_frames.py}).
      \item Compilar video MP4 con \texttt{ffmpeg}.
    \end{enumerate}
  \end{tip}
\end{titlepage}

% -------------------------------------------------------------------
\tableofcontents
\newpage

% -------------------------------------------------------------------
\section{Estructura de archivos relevantes}

\begin{center}
\begin{tabular}{lp{9cm}}
  \toprule
  \textbf{Archivo} & \textbf{Función} \\
  \midrule
  \texttt{macros/vis\_video.mac}       & Macro principal: geometría, fuente, escena, loop de eventos \\
  \texttt{macros/vis\_video\_frame.mac} & Sub-macro: un evento $+$ flush $+$ export EPS \\
  \texttt{python/accumulate\_frames.py} & Acumula PNGs frame a frame con overlay del cubo \\
  \texttt{build/}                       & Directorio de compilación; aquí se genera todo \\
  \bottomrule
\end{tabular}
\end{center}

Los comandos de esta guía se ejecutan \textbf{siempre desde \texttt{build/}},
salvo que se indique lo contrario.

% -------------------------------------------------------------------
\section{El macro principal: \texttt{vis\_video.mac}}

A continuación se muestra el macro completo con comentarios explicativos.

\lstinputlisting[style=macro, caption={macros/vis\_video.mac — macro completo}]{macros/vis_video.mac}

\subsection{Parámetros que querrás cambiar}

\begin{center}
\begin{tabular}{lll}
  \toprule
  \textbf{Parámetro} & \textbf{Línea en el macro} & \textbf{Efecto} \\
  \midrule
  Tamaño de la caja de parafina & \texttt{/detector/setParaffinX/Y/Z}
    & Dimensiones del moderador (half-extents) \\
  Distancia fuente--detector  & \texttt{/ambe/distance}
    & Posición de la fuente AmBe a lo largo de X \\
  Actividad                   & \texttt{/ambe/activity}
    & Factor de escala del espectro (neutrones/$\mu$s) \\
  Ángulo de vista             & \texttt{/vis/viewer/set/viewpointThetaPhi}
    & $\theta$ (elevación) y $\phi$ (azimut) en grados \\
  Color caja de parafina      & \texttt{/vis/geometry/set/colour Block}
    & R G B A, todos en [0,1] \\
  Color detector              & \texttt{/vis/geometry/set/colour Detector}
    & R G B A, todos en [0,1] \\
  Número de eventos           & \texttt{/control/loop \ldots{} 1 120 1}
    & Cambiar \texttt{120} al número de frames deseado \\
  \bottomrule
\end{tabular}
\end{center}

\begin{advertencia}
  \textbf{Importante — colores en Geant4:} el comando
  \texttt{/vis/geometry/set/colour} espera valores flotantes en $[0,1]$,
  \emph{no} enteros RGB de 0--255.
  Usar valores como \texttt{79 74 64} los clampea a 1.0 (blanco puro).
  Usa siempre valores como \texttt{0.47 0.43 0.35}.
\end{advertencia}

\subsection{Sub-macro de frame: \texttt{vis\_video\_frame.mac}}

Este archivo es invocado por el \texttt{/control/loop} del macro principal.
La variable \texttt{\{i\}} se sustituye automáticamente con el número de iteración.

\lstinputlisting[style=macro, caption={macros/vis\_video\_frame.mac}]{macros/vis_video_frame.mac}

Cada iteración: corre 1 evento, vuelca la escena y exporta un EPS nombrado
\texttt{frame\_\{i\}.eps} en el directorio de trabajo (\texttt{build/}).

% -------------------------------------------------------------------
\section{Paso 1 — Generar el frame de referencia del cubo}

Antes de lanzar los 120 eventos, se genera \textbf{un único frame} con la
caja de parafina sólida y sin trayectorias. Este PNG se usará después como
overlay semitransparente para que los bordes del cubo sean visibles.

\begin{lstlisting}[style=console, caption={Generar box\_frame.png desde build/}]
# Dentro de build/, con el ejecutable ya compilado:
./neutron_therm ../macros/vis_video.mac
# Geant4 abre OGLSX, configura la escena y ejecuta el loop.
# Al terminar hay 120 archivos frame_{i}.eps en build/
\end{lstlisting}

\begin{tip}
  Si sólo quieres regenerar \texttt{box\_frame.png} (por ejemplo, para cambiar
  el color de la caja sin re-correr la simulación completa), crea un macro
  mínimo:
\begin{lstlisting}[style=macro]
# box_only.mac
/detector/setParaffinX 10 cm
/detector/setParaffinY 10 cm
/detector/setParaffinZ 10 cm
/ambe/distance 10 cm
/ambe/activity 2.98
/run/initialize
/vis/open OGLSX 1280x720
/vis/viewer/set/viewpointThetaPhi 30 20
/vis/viewer/set/background 0 0 0 1
/vis/geometry/set/colour Block    0 0.65 0.60 0.48 1.0
/vis/geometry/set/colour Detector 0 1.0  0.55 0.0  1.0
/vis/viewer/set/style surface
/vis/scene/create
/vis/scene/add/volume
/vis/drawVolume
/vis/enable
/vis/viewer/refresh
/vis/ogl/set/exportFormat eps
/vis/ogl/export box_frame_0000.eps
\end{lstlisting}
y convierte con Ghostscript (ver Paso 2).
\end{tip}

% -------------------------------------------------------------------
\section{Paso 2 — Convertir EPS a PNG con Ghostscript}

Geant4 exporta en formato EPS (Encapsulated PostScript).
Ghostscript convierte cada EPS a un PNG de 1280$\times$720.

\begin{lstlisting}[style=console, caption={Conversión por lotes — desde build/}]
# Frame de referencia del cubo
gs -dNOPAUSE -dBATCH -sDEVICE=png16m -r72 -dEPSCrop \
   -sOutputFile=box_frame.png box_frame_0000.eps

# Los 120 frames de trayectorias
for i in $(seq -f "%04g" 1 120); do
    gs -dNOPAUSE -dBATCH -sDEVICE=png16m -r72 -dEPSCrop \
       -sOutputFile=frame_${i}.png frame_${i}.eps
done
\end{lstlisting}

\begin{center}
\begin{tabular}{ll}
  \toprule
  \textbf{Flag gs} & \textbf{Significado} \\
  \midrule
  \texttt{-dNOPAUSE -dBATCH} & Modo no interactivo \\
  \texttt{-sDEVICE=png16m}    & Salida PNG de 24 bits \\
  \texttt{-r72}               & 72 dpi (ajusta el tamaño de píxeles a ~1280$\times$720) \\
  \texttt{-dEPSCrop}          & Recorta al bounding box del EPS \\
  \bottomrule
\end{tabular}
\end{center}

Si quieres mayor resolución, sube \texttt{-r144} (doble) o \texttt{-r216} (triple),
y ajusta correspondientemente el tamaño de ventana OGLSX en el macro.

% -------------------------------------------------------------------
\section{Paso 3 — Acumular frames con Python}

El script \texttt{accumulate\_frames.py} genera los frames acumulados:
el frame $N$ muestra las trayectorias de los eventos 1 a $N$ superpuestas,
con el cubo de referencia como overlay tenue.

\begin{lstlisting}[style=console, caption={Acumulación — desde build/}]
source ~/mi_entorno/bin/activate

python3 ../python/accumulate_frames.py \
    --n 120 \
    --box box_frame.png \
    --box-alpha 0.25
\end{lstlisting}

\begin{center}
\begin{tabular}{lp{8.5cm}}
  \toprule
  \textbf{Argumento} & \textbf{Descripción} \\
  \midrule
  \texttt{--n 120}             & Número total de frames (debe coincidir con el loop del macro) \\
  \texttt{--box box\_frame.png} & PNG del cubo sólido de referencia \\
  \texttt{--box-alpha 0.25}    & Opacidad del overlay del cubo (0 = invisible, 1 = opaco) \\
  \texttt{--indir .}           & Directorio donde están los \texttt{frame\_NNNN.png} (default: \texttt{.}) \\
  \texttt{--outdir .}          & Directorio de salida para los \texttt{accum\_NNNN.png} (default: \texttt{.}) \\
  \bottomrule
\end{tabular}
\end{center}

Esto genera \texttt{accum\_0001.png} \ldots{} \texttt{accum\_0120.png}.

\subsection{Cómo funciona la acumulación}

\begin{enumerate}
  \item Para cada frame $N$: se calcula \texttt{np.maximum(acumulado, frame\_N)},
        manteniendo el pixel más brillante de todos los eventos anteriores.
  \item Se aplica el overlay: \texttt{np.maximum(acumulado, cubo $\times$ alpha)}.
        Como las trayectorias son blancas/amarillas (valor 255) y el overlay
        es tenue ($\sim$40), el \texttt{maximum} siempre elige la trayectoria
        donde hay una; el cubo sólo se ve donde \emph{no} hay trayectorias.
  \item El resultado se guarda como \texttt{accum\_NNNN.png}.
\end{enumerate}

\begin{tip}
  Para cambiar sólo la opacidad del cubo sin re-correr la simulación,
  basta con ajustar \texttt{--box-alpha} y re-ejecutar este paso.
  Los \texttt{frame\_NNNN.png} ya están en disco.
\end{tip}

% -------------------------------------------------------------------
\section{Paso 4 — Compilar el video con ffmpeg}

\begin{lstlisting}[style=console, caption={Compilar video MP4 — desde build/}]
ffmpeg -y \
    -framerate 60 \
    -start_number 1 \
    -i accum_%04d.png \
    -c:v libx264 \
    -pix_fmt yuv420p \
    -crf 18 \
    video_simulacion.mp4
\end{lstlisting}

\begin{center}
\begin{tabular}{lp{8.5cm}}
  \toprule
  \textbf{Flag ffmpeg} & \textbf{Descripción} \\
  \midrule
  \texttt{-framerate 60}   & 60 fps $\Rightarrow$ 120 eventos = 2 segundos de video \\
  \texttt{-start\_number 1} & El primer archivo es \texttt{accum\_0001.png} \\
  \texttt{-c:v libx264}    & Codec H.264 (compatible con casi todo) \\
  \texttt{-pix\_fmt yuv420p} & Formato de color compatible con reproductores \\
  \texttt{-crf 18}         & Calidad alta (0 = lossless, 51 = mínima) \\
  \bottomrule
\end{tabular}
\end{center}

Para cambiar la duración del video, ajusta \texttt{-framerate}:
a 30 fps serían 4 segundos; a 15 fps, 8 segundos.

% -------------------------------------------------------------------
\section{Paso 5 — Limpieza de archivos intermedios}

Los archivos EPS y los PNGs individuales pueden ocupar varios cientos de MB.
Una vez compilado el video, se pueden eliminar:

\begin{lstlisting}[style=console, caption={Limpieza — desde build/}]
rm -f frame_*.eps frame_*.png accum_*.png box_frame*.eps

# Verificar que el video quedó bien
ls -lh video_simulacion.mp4
\end{lstlisting}

% -------------------------------------------------------------------
\section{Flujo completo en un solo bloque}

El siguiente bloque ejecuta todo el pipeline de una vez, asumiendo que
ya se editó \texttt{vis\_video.mac} y que el ejecutable está compilado en \texttt{build/}.

\begin{lstlisting}[style=console, caption={Pipeline completo — desde build/}]
# 1. Simulación Geant4 (genera frame_1.eps ... frame_120.eps y box_frame_0000.eps)
./neutron_therm ../macros/vis_video.mac

# 2. Conversión EPS -> PNG
gs -dNOPAUSE -dBATCH -sDEVICE=png16m -r72 -dEPSCrop \
   -sOutputFile=box_frame.png box_frame_0000.eps

for i in $(seq -f "%04g" 1 120); do
    gs -dNOPAUSE -dBATCH -sDEVICE=png16m -r72 -dEPSCrop \
       -sOutputFile=frame_${i}.png frame_${i}.eps
done

# 3. Acumulación con overlay del cubo
source ~/mi_entorno/bin/activate
python3 ../python/accumulate_frames.py \
    --n 120 --box box_frame.png --box-alpha 0.25

# 4. Compilar video
ffmpeg -y -framerate 60 -start_number 1 -i accum_%04d.png \
    -c:v libx264 -pix_fmt yuv420p -crf 18 video_simulacion.mp4

# 5. Limpiar intermedios
rm -f frame_*.eps frame_*.png accum_*.png box_frame*.eps

echo "Listo: video_simulacion.mp4"
\end{lstlisting}

% -------------------------------------------------------------------
\section{Solución de problemas comunes}

\begin{center}
\begin{tabular}{p{4.5cm}p{9.5cm}}
  \toprule
  \textbf{Síntoma} & \textbf{Causa y solución} \\
  \midrule
  Parafina siempre blanca aunque cambies el color
    & Los valores RGB no están en $[0,1]$. Geant4 clampea integers al máximo.
      Usa \texttt{0.47 0.43 0.35}, no \texttt{79 74 64}. \\
  \addlinespace
  Los bordes del cubo salen azules en modo wireframe
    & Limitación de \texttt{gl2ps 1.4.2}: ignora el color asignado a wireframes
      y los renderiza en azul. La solución adoptada es el overlay en Python. \\
  \addlinespace
  Trayectorias tapadas por la cara frontal de la caja
    & Geant4 usa el algoritmo del pintor; la cara frontal se dibuja encima.
      Solución: modo \texttt{wireframe} en el macro + overlay en Python. \\
  \addlinespace
  Los PNG salen en blanco o muy pequeños
    & Ajusta \texttt{-r72} en Ghostscript. Prueba \texttt{-r96} si la imagen
      queda pequeña. \\
  \addlinespace
  El video dura muy poco o va muy rápido
    & Baja el framerate: \texttt{-framerate 30} (4 s) o \texttt{-framerate 15} (8 s). \\
  \addlinespace
  \texttt{./neutron\_therm: command not found}
    & Falta compilar. Desde \texttt{build/}: \texttt{cmake .. \&\& make -j4}. \\
  \bottomrule
\end{tabular}
\end{center}

% -------------------------------------------------------------------
\section{Referencia rápida de colores de partículas}

Los colores de las trayectorias se definen en \texttt{vis\_video.mac}:

\begin{lstlisting}[style=macro]
/vis/modeling/trajectories/create/drawByParticleID
/vis/modeling/trajectories/drawByParticleID-0/set neutron white
/vis/modeling/trajectories/drawByParticleID-0/set gamma  yellow
/vis/modeling/trajectories/drawByParticleID-0/set e-     red
\end{lstlisting}

Para agregar otras partículas (p.ej., protones de retroceso):
\begin{lstlisting}[style=macro]
/vis/modeling/trajectories/drawByParticleID-0/set proton cyan
\end{lstlisting}

Los nombres de colores aceptados por Geant4 incluyen:
\texttt{white}, \texttt{black}, \texttt{red}, \texttt{green}, \texttt{blue},
\texttt{yellow}, \texttt{cyan}, \texttt{magenta}, \texttt{grey}, \texttt{orange}.

\end{document}
